%=====================================
% Khai báo nhóm Tex (cơ bản)
%=====================================
%===========Bảng
\usepackage{longtable,multirow,makecell}
\usepackage[dvipsnames,table]{xcolor}
\usepackage{color,colortbl}
\usepackage{pgf-pie,diagbox}
\usepackage{multirow}
\usepackage{setspace}
\usepackage{ragged2e}
\usepackage[utf8]{vietnam}
\usepackage[version=3]{mhchem}
%---
\usepackage{extarrows,amsmath,amssymb,mathrsfs,maybemath,xlop,polynom,slashbox}
%\usepackage{mathptmx}
\usepackage{yhmath} %\let\widering\relax
\usepackage{enumerate}
\usepackage{enumitem} % định dạng môi trường liệt kê 
\setlist[itemize,enumerate,description]{noitemsep, {nosep}}
\usepackage{tikz} 
\usepackage{tkz-euclide}
%\usepackage{ex_tkz-euclide}
\usepackage{tikz-3dplot}
\usepackage{tkz-tab}
\usepackage{pifont}
\usepackage{bbding}
\usepackage{stackengine}
\usepackage{scalerel}
\usepackage{array}
\usepackage{xcolor,colortbl,tasks}
%==========
\usetikzlibrary{math,through,calc,intersections,angles,quotes,shapes,shapes.geometric,arrows,patterns,snakes,matrix,chains,arrows.meta,decorations.shapes,decorations.fractals,decorations.markings,shadows}
\usetikzlibrary{positioning,decorations.text,decorations.pathmorphing}% Để uốn cong văn bản 
\usetikzlibrary{shadings,fadings} %ĐỔ BÓNG
\usepackage{pgfplots}
\usepackage{pgfornament}
\usepgfplotslibrary{fillbetween}
\pgfplotsset{compat=1.13}
\usepackage[hidelinks,unicode]{hyperref}
\usepackage{currfile}
\usepackage[outline]{contour} %viền
\usepackage{fontawesome} % Gói kí hiệu
\usepackage{lipsum} %Lấy text
\usepackage{capt-of}
\usepackage[locale=DE]{siunitx} % cách viết số đo có đơn vị theo chuẩn DE (gần giống VN) 
%===========Bảng
%\usepackage{longtable,multirow,makecell}
%\usepackage{diagbox}
%\renewcommand{\tabcolsep}{5mm}
\arrayrulecolor{teal}
\newcolumntype{C}[1]{>{\centering\arraybackslash}p{#1}}
\newcolumntype{L}[1]{>{\raggedright\arraybackslash}p{#1}}
%----------
\pgfmathsetmacro{\mepphai}{\phai+\leftnote} 
\usepackage[top=\tren cm, bottom=\duoi cm, left=\trai cm, right=\mepphai cm] {geometry}
\usepackage{varwidth}
%--------------Gói trắc nghiệm EX-TEST
\usepackage[loigiai]{ex_test} 
%\usepackage[solcolor]{ex_test}
%%--------Định nghĩa lại hiển thị Lời giải, đn lại chữ CÂU
\font\damEX=ugqb8v at 11.5pt
%\def\loigiaiEX{
%	\vspace*{-4mm}
%	\begin{center}
%		\color{\mauLG}\fontfamily{qag}\bfseries\strut%\faFolderOpen\ Lời giải.
%		\faCommenting\ Lời giải.
%	\end{center}
%	\vspace*{-3mm}
%}
%--------Định nghĩa lại hiển thị Lời giải, đn lại chữ CÂU
\def\loigiaiEX{\damEX\color{\mauLG}\faCommenting\ Lời giải.%
}
%----Viết tắt icon cho VD, BT, EX; Hiền thị tên EX; Dấu kết thúc
%\def\iconVD{\faGg\,}
\renewcommand{\nameex}{\damEX\color{\mauEX} \iconVD Câu}
\def\qedEX{\color{\mauDC}\ensuremath{\square}}
%---------- Khai báo viết tắt, in đáp án
\newcommand{\hoac}[1]{ %hệ hoặc
	\left[\begin{aligned}#1\end{aligned}\right.}
\newcommand{\heva}[1]{ %hệ và
	\left\{\begin{aligned}#1\end{aligned}\right.}
%--In đáp án
\newcommand{\indapan}[2]{
	\addcontentsline{toc}{subsection}{\sf Bảng đáp án} % đưa MT vào mục lục
	\begin{center}
		\begin{tikzpicture}%
			\node[thick,scale=1,fill=\mausubsec!3,draw=\mausubsec,minimum width=3.5cm,minimum height=0.1cm,rounded corners=2mm]{\fontfamily{qag}\fontsize{11}{11}\selectfont\bfseries\color{\mausubsec} BẢNG ĐÁP ÁN};
		\end{tikzpicture}%
	\end{center}
	\inputansbox{#1}{#2}
}
%----------
\usepackage{esvect}
\def\vec{\vv} %vecto
\def\overrightarrow{\vv}
%Lệnh song song
\DeclareSymbolFont{symbolsC}{U}{txsyc}{m}{n}
\DeclareMathSymbol{\varparallel}{\mathrel}{symbolsC}{9}
\DeclareMathSymbol{\parallel}{\mathrel}{symbolsC}{9}
%--------------------------
% HEADER AND FOOTER STYLING
%--------------------------
%Tạo background
\def\hienLOGO{
	\usepackage{background}
	\backgroundsetup{%
		scale=1,%%
		angle=0,%%
		opacity=0.2,%%
		contents={
			\ifnum\the\value{page}>0%
			\begin{tikzpicture}[remember picture,overlay,scale=1]
				\node at (current page.center){\includegraphics[width=13cm]{logo/logo.jpg}};
			\end{tikzpicture}
			\else
			\fi
		}
	}	
}
%----
\def\anLOGO{}
%=======================
\newcommand{\myfancyhead}{% trên và chấm trái
	\boldmath
	\begin{tikzpicture}[remember picture,overlay,>=stealth]
		\fill[fill=\maufoot] ([yshift=-\tren cm+\topset cm]current page.north west) rectangle +(\paperwidth,0.5mm) coordinate (AA);
		\foreach \ii/\mauphu in {5/white,4.8/\maufootmot,4/white,3.8/\maufoothai,3/white,2.8/\maufootba,2/white,1.8/\maufoot}{
			\fill[fill=\mauphu,draw=white] ([yshift=-2mm-0.5*\topset cm]AA)--++(-\ii,0)--++(0.5*\topset cm+2mm,0.5*\topset cm+ 2mm)--++(\ii cm,0)--cycle;}
		%----
		\checkoddpage\ifoddpage %nếu trang lẻ
		\def\kctrai{\trai} \def\kcphai{\phai} \def\ndtren{\leftmark}
		\else
		\def\kctrai{\phai} \def\kcphai{\trai} \def\ndtren{\rightmark}
		\fi
		%------
		\node[anchor=south west,text=\mauchufoot,inner sep=0pt] at ([yshift=-\tren cm+\topset cm+1mm,xshift=\kctrai cm]current page.north west) {\fontfamily{qag}\fontsize{8pt}{1pt}\selectfont \nouppercase{\ndtren}};
		\draw ([yshift=-\tren cm+\topset cm+1mm,xshift=-\kcphai cm]current page.north east) node[anchor=south east,text=\mauchufoot,inner sep=0pt] {\fontfamily{qag}\fontsize{8pt}{1pt}\selectfont \trenphai};
		%--tròn trái
		\fill[\maufoot] (current page.north west) circle (\tren -\topset);
		\draw[white,line width=1pt] (current page.north west) circle (\tren cm -\topset cm-3pt);
		%--Trang
		\draw (current page.north west)+(-45:0.5*\tren -0.5*\topset) node[white]{\fontfamily{qag}\fontsize{7pt}{1pt}\selectfont\bfseries\thepage};
		%=====================Dòng chấm note
		\path ([yshift=-\tren cm+\topset cm]current page.north west) coordinate (AA)
		++(\paperwidth,0)coordinate (BB);
		\checkoddpage\ifoddpage %nếu trang lẻ
		%----- Kẻ đứng
		\draw[\maufoot] ([xshift=-\mepphai cm+2.5mm]BB)--([yshift=\duoi cm-0.75*\botset cm+1mm,xshift=-\mepphai cm+2.5mm]current page.south east);
		%--		
		\path ([yshift=-\tren cm+0.5*\topset cm-0.5cm,xshift=-\phai cm-0.5*\leftnote cm+2.5mm]current page.north east) coordinate (DDD); 
		\begin{scope}
			\clip ([yshift=-\tren cm+0.5*\topset cm-1pt,xshift=-\phai cm]current page.north east) rectangle ([yshift=\duoi cm-0.5*\botset cm,xshift=-\mepphai cm+5mm]current page.south east);% cắt chấm
			\node[inner sep =0pt,scale=1,anchor=north] at ([yshift=0cm,xshift=0pt]DDD) {
				\parbox{\leftnote cm}{\centering
					\dotlineEX{60}
				}
			};
		\end{scope}
		\else %chẵn
		%----- Kẻ đứng
		\draw[\maufoot] ([xshift=\mepphai cm-2.5mm]AA)--([yshift=\duoi cm-0.75*\botset cm+1mm,xshift=\mepphai cm-2.5mm]current page.south west);
		%--		
		\path ([yshift=-\tren cm+0.5*\topset cm-0.5cm,xshift=\phai cm+0.5*\leftnote cm-2.5mm]current page.north west) coordinate (DDD); 
		\begin{scope}
			\clip ([yshift=-\tren cm+0.5*\topset cm-1pt,xshift=\phai cm]current page.north west) rectangle ([yshift=\duoi cm-0.5*\botset cm,xshift=\mepphai cm-5mm]current page.south west);% cắt chấm
			\node[inner sep =0pt,scale=1,anchor=north] at ([yshift=0cm,xshift=0pt]DDD) {
				\parbox{\leftnote cm}{\centering
					\dotlineEX{60}
				}
			};
		\end{scope}					
		\fi
		%===================Logo và Quick note
		\checkoddpage\ifoddpage %nếu trang lẻ
		%--tiêu đề phải
		\path ([yshift=-\tren cm+0.5*\topset cm-0.5cm,xshift=-\phai cm-0.5*\leftnote cm+2.5mm]current page.north east) coordinate (DD); 
		%--
		\fill[white] ([yshift=-\tren cm+0.5*\topset cm-5pt,xshift=\trai cm-2pt]current page.north east) rectangle ([yshift=\duoi cm-0.5*\botset cm-8cm,xshift=-\mepphai cm+3mm]current page.north east);
		\node[inner sep =0pt,anchor=north] (thanhta) at ([yshift=-2mm]DD) {
			\includegraphics[width=\leftnote cm-0.5cm]{logo/logo.jpg}		
		};	
		\else
		%--tiêu đề phải
		\path ([yshift=-\tren cm+0.5*\topset cm-0.5cm,xshift=\phai cm+0.5*\leftnote cm-2.5mm]current page.north west) coordinate (DD); 
		%--
		\fill[white] ([yshift=-\tren cm+0.5*\topset cm-5pt,xshift=\phai cm-2pt]current page.north west) rectangle ([yshift=\duoi cm-0.5*\botset cm-8cm,xshift=\mepphai cm-3mm]current page.north west);
		\node[inner sep =0pt,anchor=north] (thanhta) at ([yshift=-2mm]DD) {
			\includegraphics[width=\leftnote cm-0.5cm]{logo/logo.jpg}		
		};
		\fi
		%\draw (thanhta.south) node[below=0pt,xscale=0.8]{\small\normalfont\color{\mauname} Sưu tầm \& Biên tập};
		%---note
		\node[inner sep =6pt, text=black,scale=1,anchor=north,fill=\maufoot!3,draw=\maufoot] (bon) at ([yshift=-1cm]thanhta.south) {
			\parbox{\leftnote cm-5mm-12pt}{ \fontsize{10}{15}\selectfont\normalfont
				\vspace*{2pt}
				\chamngon%
			}
		};
		\draw[\maufoot, line width=5pt] (bon.north west)--(bon.north east);
		\draw[\maufoot!50] ([yshift=9pt,line width=0.4pt]bon.north east)--([yshift=9pt]bon.north west)
		node[fill=white,inner sep=2pt,anchor=south west,yshift=-2pt,xshift=-2pt]{\bfseries\color{\maufoot}ĐIỂM:}
		;
		%--note dưới
		\node[inner sep =6pt, text=white,scale=1,anchor=north,fill=\maufoot] (noteduoi) at ([yshift=-0.25cm]bon.south) {
			\parbox{\leftnote cm-5mm-12pt}{\fontfamily{qag}\fontsize{11}{1}\selectfont\bfseries\centering
				QUICK NOTE
			}
		};
		\draw[\maufoot, line width=0.4pt] ([yshift=-2pt]noteduoi.south west)--([yshift=-2pt]noteduoi.south east);
	\end{tikzpicture}%
}
% trên mục lục
\newcommand{\headmucluc}{%
	\boldmath
	\begin{tikzpicture}[remember picture,overlay,>=stealth]
		\fill[fill=\maufoot] ([yshift=-\tren cm+\topset cm]current page.north west) rectangle +(\paperwidth,0.5mm) coordinate (AA);
		\foreach \ii/\mauphu in {5/white,4.8/\maufootmot,4/white,3.8/\maufoothai,3/white,2.8/\maufootba,2/white,1.8/\maufoot}{
			\fill[fill=\mauphu,draw=white] ([yshift=-2mm-0.5*\topset cm]AA)--++(-\ii,0)--++(0.5*\topset cm+2mm,0.5*\topset cm+ 2mm)--++(\ii cm,0)--cycle;}
		%----
		\checkoddpage\ifoddpage %nếu trang lẻ
		\def\kctrai{\trai} \def\kcphai{\phai} \def\ndtren{\leftmark}
		\else
		\def\kctrai{\phai} \def\kcphai{\trai} \def\ndtren{\rightmark}
		\fi
		%------
		\node[anchor=south west,text=\mauchufoot,inner sep=0pt] at ([yshift=-\tren cm+\topset cm+1mm,xshift=\kctrai cm]current page.north west) {\fontfamily{qag}\fontsize{8pt}{1pt}\selectfont \nouppercase{\ndtren}};
		\draw ([yshift=-\tren cm+\topset cm+1mm,xshift=-\kcphai cm]current page.north east) node[anchor=south east,text=\mauchufoot,inner sep=0pt] {\fontfamily{qag}\fontsize{8pt}{1pt}\selectfont \trenphai};
		%--tròn trái
		\fill[\maufoot] (current page.north west) circle (\tren -\topset);
		\draw[white,line width=1pt] (current page.north west) circle (\tren cm -\topset cm-3pt);
		%--Trang
		\draw (current page.north west)+(-45:0.5*\tren -0.5*\topset) node[white]{\fontfamily{qag}\fontsize{7pt}{1pt}\selectfont\bfseries\thepage};
	\end{tikzpicture}%
}
%===========================
\newcommand{\myfancyfoot}{% dưới
	\begin{tikzpicture}[remember picture,overlay]
		\coordinate (AA) at ([yshift=\duoi cm-0.75*\topset cm-0.5mm]current page.south west);
		\checkoddpage\ifoddpage %nếu trang lẻ
		\def\kctrai{\trai} \def\kcphai{\phai}
		\else
		\def\kctrai{\phai} \def\kcphai{\trai}
		\fi
		%------
		%--Đặt tạm để đo độ dài
		\node[anchor=east,text=white,inner sep=0pt] (sotrang) at ([yshift=-0.3cm]AA)  {\fontfamily{qag}\fontsize{8pt}{1pt}\selectfont{\bfseries \thepage}{\footnotesize /}{\bfseries\pageref{LastPage}}};
		%---Đo độ dài chữ chương
		\path (sotrang.west) coordinate(chA) (sotrang.east) coordinate(chB);
		\path let \p1 =($(chA)-(chB)$),\n1= {scalar(veclen(\x1,\y1)/1cm)}in \pgfextra{\xdef\kcthem{\n1}};
		%--
		\pgfmathsetmacro{\daisach}{\kctrai+\kcthem-0.1} 
		\fill[\maufoot] ([xshift=\daisach cm+10mm]AA) rectangle +(\paperwidth,0.5mm) coordinate (BB);
		%--
		\fill[fill=\maufoot!20] ([yshift=1mm]AA)--++(\daisach cm,0) to[bend left=8mm] ++(3mm,-1.5mm)coordinate (Tr1)--++(0,-0.75)to[bend right=8mm]++(-3mm,1.5mm)--++(-\daisach,0)--cycle;
		\fill[fill=\maufoot] (Tr1) to[bend left=8mm]++(0.6,3mm)--++(0,-0.75)to[bend right=8mm]([yshift=-0.75cm]Tr1)--cycle;
		\draw[\maufoot,line width=0.5mm] ([yshift=1mm]AA)--++(\daisach cm,0) to[bend left=8mm] ++(3mm,-1.5mm);
		%
		\node[anchor=east,text=\mauchufoot,inner sep=0pt] (sotrang) at ([yshift=-0.25cm,xshift=-5pt]Tr1)  {\fontfamily{qag}\fontsize{8pt}{1pt}\selectfont{\bfseries \thepage}{\footnotesize /}{\bfseries\pageref{LastPage}}};
		%		\draw ([yshift=-0.25cm,xshift=-3pt]Tr1) node[anchor=east,text=\mauchufoot,inner sep=0pt] {\fontfamily{qag}\fontsize{8pt}{1pt}\selectfont{\bfseries \thepage}{\footnotesize /}{\bfseries\pageref{LastPage}}};
		%
		\node[anchor=north east,text=\mauchufoot,inner sep=0pt] (thinh) at ($([xshift=-\kcphai cm]AA)+(\paperwidth,-2pt)$) {\fontfamily{qag}\fontsize{10pt}{1pt}\selectfont%
			\duoiphai%
%			\faFileText\, \hoten \,\,--\,\, \faPhoneSquare\ \SDT%
			%			\faFileText\, Tài liệu Toán 11 - Chương trình mới%
		};
		%--tròn phải
		\fill[\maufoot] (current page.south east) circle (\duoi -0.75*\topset);
		\draw[white,line width=1.5pt] (current page.south east) circle (\duoi cm -0.75*\topset cm-3pt);
	\end{tikzpicture}%
}
%======================Head chapter theo note, fullwidth
\usepackage{changepage}
\strictpagecheck
\usepackage{lastpage}
\usepackage{fancyhdr,lastpage}
\pagestyle{fancy}
\fancyhf{}
\def\headchuong{\myfancyhead}
\fancypagestyle{plain}{
	\fancyhead[LO,RE]{\headchuong}
	\fancyfoot[LO,RE]{\myfancyfoot}
}
\fancyhead[LO,RE]{\myfancyhead}
\fancyfoot[LO,RE]{\myfancyfoot}
\renewcommand{\footrulewidth}{0pt}
\renewcommand{\headrulewidth}{0pt}
%--------------4.2
\usepackage[most]{tcolorbox}
\colorlet{tcbcol@back}{tcbcolback}
\colorlet{tcbcol@frame}{tcbcolframe}
%-------------- Khung (Trong main này ko sd)
\newtcolorbox[auto counter]{khung4}[1]{enhanced jigsaw, breakable,
	before skip=1mm,after skip=1mm,
	left=1mm,right=1mm,top=2mm,bottom=1mm,
	colframe=cyan!70!black,colback=yellow!15,colbacktitle=cyan!10,coltitle=red!70!black,
	opacityback=0.5,
	,boxrule=0.4mm,
	attach boxed title to top center=
	{yshift=-0.1mm-\tcboxedtitleheight/2,yshifttext=2mm-\tcboxedtitleheight/2},
	varwidth boxed title*=-3cm,
	boxed title style={boxrule=0.3mm,
		frame code={ \path[tcb fill frame] ([xshift=-4mm]frame.west)
			-- (frame.north west) -- (frame.north east) -- ([xshift=4mm]frame.east)
			-- (frame.south east) -- (frame.south west) -- cycle; },
		interior code={ \path[tcb fill interior] ([xshift=-2mm]interior.west)
			-- (interior.north west) -- (interior.north east)
			-- ([xshift=2mm]interior.east) -- (interior.south east) -- (interior.south west)
			-- cycle;} 
	},
	fonttitle=\fontfamily{qag}\fontsize{10}{0}
	\bfseries,
	fontupper=\sffamily,
	title={#1}
}
%---------------Dạng toán
\newenvironment{tttt}{}{}
\newcounter{dang}\setcounter{dang}{0}
\renewcommand{\thedang}{\arabic{dang}}
%---Dạng 1
\newtcolorbox{dangg}[1]{%
	fonttitle=\fontfamily{qag}\bfseries,%fontupper=\itshape,
	enhanced jigsaw,
	colframe=\maudang,colback=\maunendang,coltitle=white,%
	sharp corners, breakable, halign title=center,%adjusted title=center, %canh giữa DẠNG
	opacityback=0.5,
	before skip=2mm,after skip=3mm,%
	left=2mm,right=2mm,top=2mm,bottom=2mm,%
	boxrule=1pt,%
	title={\faFolderOpen\ Dạng~\stepcounter{dang}\thedang.\ #1}%
%	\setcounter{subsubsection}{0}%
	\addcontentsline{toc}{subsection}{\fontsize{11}{11}\selectfont\color{\mauSUBSUBSEC}\it\sffamily \faFolderOpen\ Dạng~\thedang.~#1~}%
}
%====
\newenvironment{dang}[1]{
	\begin{dangg}{#1}
	}
	{\setcounter{vd}{0}
	\end{dangg}
}

%\newcommand{\dang}[2]{
%\begin{dangg}{#1}
%	\begin{tomtat}
%	
%	\end{tomtat}
%\end{dangg}
%}
%---------------------------------------------------------------
% ĐỊNH NGHĨA SECTION. SUBSECTION, SUBSUBSECTION ... THEO Ý RIÊNG
%---------------------------------------------------------------
\usepackage[explicit]{titlesec} % để gọi #1
\usepackage{titledot} % gói lệnh chứa cả titlesec và titletoc
%=====================================
\setcounter{secnumdepth}{4} %độ sâu
\renewcommand{\thechapter}{\arabic{chapter}}
\renewcommand{\thesection}{\arabic{section}}
\renewcommand{\thesubsection}{\Alph{subsection}}
\renewcommand{\thesubsubsection}{\arabic{subsubsection}}
%--------------Tròn
\newcommand{\tron}[1]{
	\begin{tikzpicture}[baseline=(A.base)]%
		\node[circle,draw=\mauSUBSEC,line width=0.5pt,fill=white,inner sep=2pt,outer sep=1.5pt] (A) {\color{\mauSUBSEC} #1};
		\node[circle,draw=white,fill=\mauSUBSEC,inner sep=1pt,outer sep=1pt] (A) {\bfseries\color{white} #1};
	\end{tikzpicture}%
}
%============================Mục lục - Chapter*
%\makeatletter
\titleformat{name=\chapter,numberless}[display]
{\color{\mauchuong}\fontsize{23pt}{1pt}\fontfamily{qag}\selectfont\bfseries}
{}
{-1em}
{%
	\begin{tcolorbox}		
		[
		enhanced,
		frame empty,
		top=0.2cm, 
		bottom=0.2cm,
		colback=white,
		overlay={
			\fill[\mauchuong] (frame.south west) rectangle ([yshift=-0.5cm]frame.north east);
			\fill[\maunenmucluc,rounded corners=5pt,draw=black,line width=1pt] (frame.south west)--(frame.south east)--++(0,0.25)--++(-6,0)--([xshift=2cm]frame.north)--++(-4,0)--([xshift=6cm,yshift=0.25cm]frame.south west)--([yshift=0.25cm]frame.south west)--cycle;
			%---
			\fill[\mauchuong,opacity=0.25] (frame.south west) rectangle ([yshift=-0.75cm]frame.south east);
			\fill[\maunenmucluc,rounded corners=5pt,draw=black,line width=1pt,opacity=0.25] (frame.south west)--(frame.south east)--++(0,-0.25)--++(-6,0)--([xshift=2cm,yshift=-1cm]frame.south)--++(-4,0)--([xshift=6cm,yshift=-0.25cm]frame.south west)--([yshift=-0.25cm]frame.south west)--cycle; 
			%--- 
			\foreach \iii/\mau/\bk in {-2/\maumlba/0.15,-1/\maumlhai/0.2,0/\maumlmot/0.25,1/\maumlhai/0.2,2/\maumlba/0.15}{
				\fill[draw=black,ball color=\mau]([yshift=-0.4cm,xshift=0.7*\iii cm]frame.south) circle (\bk cm);
			}
			\node[xshift=0mm, yshift=0cm,scale=1,right=0cm] (tentron) at ([yshift=-0.4cm,xshift=0.7*2.1 cm]frame.south) {\pgfornament[height=0.5cm,color=\maumlmot,opacity=1]{14}};
			\node[xshift=0mm, yshift=0cm,scale=1,left=0cm] (tentron) at ([yshift=-0.4cm,xshift=-0.7*2.1 cm]frame.south) {\pgfornament[height=0.5cm,color=\maumlmot,opacity=1]{11}};
		}
		]
		\color{\mauchuong}\fontsize{23pt}{1pt}\fontfamily{qag}\selectfont\bfseries\centering\MakeUppercase{#1}
	\end{tcolorbox}%
	\begin{tikzpicture}[overlay,remember picture]
		\foreach \iii in {0,1,...,22}{\draw[\mauchuong!50,line width=2pt,opacity=0.2] 
			([xshift=0.3*\iii cm]current page.north west)--++(-135:10cm)
			;}
	\end{tikzpicture}
}
[\vspace{1cm}]
%-------Đn section---------------------------
%--Biến đếm cho các tùy chọn BT
\newcounter{stt}
\setcounter{stt}{0}
%--
\titlespacing*{\section}{0cm}{0cm}{0cm}[0cm]
\titleformat
{\section}
{\color{white}\fontfamily{put}\fontsize{16pt}{22pt}\selectfont\bfseries}
{}
{0cm}
{%	
	\setcounter{ex}{0}%
	\setcounter{vd}{0}%
	\setcounter{bt}{0}%
	\setcounter{dang}{0}%
	\setcounter{stt}{0}%
	%	\setcounter{vidu}{0}%
	\begin{tikzpicture}[line join=round,line cap=round]
		\def\canhlg{0.65}
		\pgfmathsetmacro{\kchm}{2*\canhlg*cos(30)} 
		\def\ndbai{\parbox{\textwidth-\kchm cm-2cm}{\centering{\MakeUppercase{#1}}}}
		%-- Lục giác
		\fill[\maubai] (0,0)--++(-30:\canhlg)--++(30:\canhlg cm)--++(90:\canhlg cm)--++(150:\canhlg cm)--++(210:\canhlg cm)--cycle; 		
		%-- kẻ dưới
		\draw[draw=\mauchuong, line width=2pt] 
		(0,0)++(-30:\canhlg)++(0,-5pt) --++(30:\canhlg)--++(\textwidth-\kchm cm-1.3 cm,0) coordinate (tronphai)
		(0,0)++(-30:\canhlg)++(0,-5pt)--++(150:\canhlg)--++(-1.1,0) coordinate (bai);
		\fill[fill=\mauchuong] (tronphai) circle (3pt) (bai) circle (3pt);
		%-- Nội dung chính --	
		\draw (0,0)++(-30:\canhlg)++(0,-5pt) ++(30:\canhlg)++(0.5*\textwidth-0.5*\kchm cm-0.7 cm,0) node[inner sep=0pt,text=\maubai,above=3pt] {\ndbai};
		%--Bài và số bài
		\node[text=white] at (30:\canhlg) {\usefont{T5}{HLButlong}{m}{n}\fontsize{20pt}{1pt}\selectfont\thesection};
		\draw ([yshift=4mm,xshift=-3pt]bai) node[inner sep=0pt,right,\maubai,scale=0.95] {\usefont{T5}{HLButlong}{m}{n}Bài};
\end{tikzpicture}}
[]
%-------Đn section***---------------------------
\titleformat
{name=\section,numberless}
{\color{white}\fontfamily{put}\fontsize{16pt}{22pt}\selectfont\bfseries}
{}
{0mm}
{%	
	\setcounter{ex}{0}%
	\setcounter{vd}{0}%
	\setcounter{bt}{0}%
	\setcounter{dang}{0}%
	\setcounter{stt}{0}%
	%	\setcounter{vidu}{0}%
	\setcounter{subsection}{0}%
	\setcounter{subsubsection}{0}%
	\addcontentsline{toc}{section}{\hspace*{-3.8pc}\fontsize{12.5pt}{13pt}\selectfont\sffamily\color{\mauSEC} #1}%
	\begin{tikzpicture}[line join=round,line cap=round]
		\def\ndbai{\parbox{\textwidth}{\centering{\MakeUppercase{#1}}}}		
		%-- kẻ dưới
		\draw[draw=\mauchuong, line width=2pt] 
		(0,0)--++(\textwidth-5pt,0) coordinate (tronphai);
		\fill[fill=\mauchuong] (tronphai) circle (3pt) (0,0) circle (3pt);
		%-- Nội dung chính --	
		\draw (0,0)++(0.5*\textwidth,0) node[inner sep=0pt,text=\maubai,above=2pt] {\ndbai};
\end{tikzpicture}}
[]
%-------Đn subsection---------------------------
%\newcommand{\dnsub}[1]{%
%	\setcounter{ex}{0}%
%	\setcounter{vd}{0}%
%	\setcounter{bt}{0}%
%	\begin{tikzpicture}[line join=round,line cap=round]%
%		\node[inner sep=3pt,text=white,rounded corners=0pt,fill=\mausubsec] (muccon) at (0,0){\hspace*{2mm}#1\hspace*{0.75cm}}; %trái
%		\fill[fill=\mausubsecphu] ([xshift=-3pt]muccon.north east) --([xshift=-3pt]muccon.east) 
%		.. controls +(-90:0.2) and +(10:0.2) ..  ([xshift=-0.35cm]muccon.south east)
%		--(muccon.south east) .. controls +(10:0.2) and +(-90:0.2) .. ([xshift=0.35cm]muccon.east)--([xshift=0.35cm]muccon.north east)--cycle ; %phải
%	\end{tikzpicture}%
%}
%%\titlespacing{\subsection}{0cm}{0cm}{0cm}[0cm]
%\titleformat{\subsection}
%{\normalfont\fontsize{13pt}{13pt}\fontfamily{qag}\selectfont\bfseries}
%{}
%{-1.5mm}
%{	
%	\dnsub{\thesubsection\,--\,\MakeUppercase{#1}}
%}
%[]
%-------Đn subsection---------------------------
\newcommand{\dnsubb}[1]{
	\begin{tikzpicture}[baseline=(muccon.base),line join=round,line cap=round]%
		\node[diamond,inner sep=4pt,text=white,fill=\mausubsecphu!50,draw=\mausubsecphu!50] at (4pt,0pt){A};
		\node[diamond,inner sep=4pt,text=white,fill=white,draw=\mausubsec] (muccon) at (0,0){#1}; %trái
		\node[diamond,inner sep=2.5pt,text=white,ball color=\mausubsec] (muccon) at (0,0){#1}; %trái
	\end{tikzpicture}%
}
\titlespacing{\subsection}{0cm}{0cm}{2mm}[0cm]
\titleformat{\subsection}
{\normalfont\fontsize{13pt}{13pt}\fontfamily{qag}\selectfont\bfseries\color{\mausubsec}}
{\hspace*{-2mm}\dnsubb{\fontsize{15}{1}\selectfont\thesubsection}}
{1mm}
{	
	\MakeUppercase{#1}
}
[]
%----------ĐN subsubsection-----------------------
%\titlespacing{\subsubsection}{0pt}{4mm}{1mm}[0cm]
\titleformat{\subsubsection}
{
	\fontsize{11.5pt}{9pt}\fontfamily{qag}\selectfont\bfseries\color{\mausubsubsec}}
{
	\hspace*{-2.5mm}
	\begin{tikzpicture}[baseline=(char.base),line join=round,line cap=round]
		\node[text=white, fill=\mausubsubsec,rounded corners=3pt] (char) at (0,0){\,\thesubsubsection\,};
	\end{tikzpicture}%
}
{0.75em}
{#1}
[]
%----------ĐN paragraph-----------------------
\titlespacing{\paragraph}{0pt}{0mm}{0mm}[0cm]
\titleformat{\paragraph}
{\fontsize{11.5pt}{17pt}\fontfamily{put}\selectfont\bfseries\color{\maupara}}
{\theparagraph.}
{0.5em}
{#1}
[]
%================= Đn chương
\font\damC=ugqb8v at 24pt
\titlespacing{\chapter}{0cm}{0.5cm}{0.25cm}[0cm] %1: , 2: Trên, 3: dưới
\titleformat{\chapter}
{\fontsize{25pt}{25pt}\fontfamily{qag}\selectfont\bfseries\color{white}} %định dạng chung
{} %đánh số
{-0.4cm}
{
	\parbox{\textwidth}{\damC%
		\raggedleft\MakeUppercase{#1}} %neo vitri
	\begin{tikzpicture}[remember picture,overlay,line join=round,line cap=round] 
		%------------------------------------
		\checkoddpage\ifoddpage %nếu trang lẻ
		\def\sokc{\mepphai}
		\else
		\def\sokc{\trai}
		\fi
		%==============
		\node[scale=2,text=\mauchuong,anchor=north east,inner sep=0pt]  (tam) at ([xshift=-\sokc cm,yshift=-1.25*\tren cm]current page.north east){
			\usefont{T5}{HLButlong}{m}{n}
			{\thechapter}
		};
		%--	chương
		\node[anchor=south east, text=\mauchuong, inner sep=0pt,scale=0.85,yshift=-2pt] (chuong) at ([xshift=-5mm,yshift=2mm]tam.south west){
			\usefont{T5}{HLButlong}{m}{n}\chaptername
		};
		%============= Nền màu trái
		\begin{scope}
			\clip (tam.north east) rectangle +(-\textwidth,-1.5cm);
			\foreach \ii/\mau in {0/\mauchuong!30,0.5/\mauchuong!70,1.2/\mauchuong}{
				\fill[\mau,draw=white,line width=4pt] 
				([xshift=-4mm-\ii cm]chuong.north west)
				--++(-0.5,-0.88)--++(-\textwidth,0)--++(0,0.85)--cycle;
			}
		\end{scope}
		%-- nội dung chương
		\node[scale=1,text=\mauchuong,anchor=north east,inner sep=0pt]  at 	([yshift=-0.25cm]tam.south east) {
			\parbox{\textwidth}{\damC%
				\raggedleft\MakeUppercase{#1}}
		};
	\end{tikzpicture}
}
[\vspace{0.5cm}]
%================== ĐN Part
\font\damCHUONG=ugqb8v at 16pt
\titlespacing{\part}{0cm}{0cm}{0cm}[0cm]
\titleformat{\part}[display]
{\bfseries}
{}
{0pt}
{
	\begin{tikzpicture}[remember picture, overlay]
		\checkoddpage\ifoddpage %nếu trang lẻ
		\def\kcphu{\phai}
		\else
		\def\kcphu{\trai}
		\fi
		\fill[\nenphan] (current page.north west) rectangle (current page.south east);
		%----Tên phần
		\node[inner sep=0pt,scale=2,left] (tenphan) at ([xshift=-\kcphu cm,yshift=-7cm]current page.north east) {\color{\mauphan!70!black}\fontfamily{qag}\bfseries\fontsize{15pt}{1pt}\selectfont\MakeUppercase{\partname}\fontsize{80pt}{1pt}\selectfont \thepart};
		%----Nền màu trái
		\foreach \ii/\mau in {0/\mauphan!30,0.5/\mauphan!70,1.2/\mauphan}{
			\fill[\mau,draw=\nenphan,line width=4pt] 
			([xshift=-4mm-\ii cm,yshift=1cm]tenphan.south west)
			--++(-0.5,-1.2)--++(-\paperwidth,0)--++(0,1.2)--cycle;
		}
		%-----Nội dung phần
		\node[inner sep=0pt,below left,scale=2.5] at ([xshift=-4mm,yshift=-0.5cm]tenphan.south east) 
		{
			\parbox{6.5cm}{
				%\fontsize{20pt}{20pt}\fontfamily{qag}\selectfont\bfseries%
				\damCHUONG\color{\mauchuphan}\raggedleft
				\MakeUppercase{#1}
			}
		};
	\end{tikzpicture}
}
[
\vspace{2cm}
\thispagestyle{empty}
]
%============================ 
%\def\mauitemCI{cyan}\def\mauitemKN{red}
\def\itemKN{\color{\mauitemKN}\faCheckSquareO}
\def\itemCI{\color{\mauitemCI}\faCheckCircleO}
%%======= Thiết lập labelitem, labelenumerate
%\renewcommand{\labelitemi}{\color{red}\faCheckSquareO}
\renewcommand{\labelitemi}{\color{\mauitemCI}\faCheckCircleO}
\renewcommand{\labelitemii}{\color{\mauitemCI}\bf ---}
\renewcommand{\labelitemiii}{\color{\mauitemCI}\bf +}
\renewcommand{\labelenumi}{\alph{enumi})}
%\renewcommand{\labelenumii}{\color{blue}\bf\arabic{enumi}.\arabic{enumii}}
%============================
%============================
% Canh chỉnh mục lục chính
\setcounter{secnumdepth}{4} %Độ sâu đánh số
\setcounter{tocdepth}{1} %Độ sâu mục lục
\contentsmargin{0cm}
%%~~~~~~~~~~~~~~~~~~~~~
%\def\mauPART{red!70!black}
%\def\mauCHAP{red}
%\def\mauSEC{magenta}
%\def\mauSUBSEC{blue}
%\def\mauSUBSUBSEC{blue!80}
%----
\titlecontents{part}[0pc]
{\addvspace{10pt}%
	\color{\mauPART}\fontsize{18pt}{1pt}\fontfamily{qag}\selectfont\bfseries 
}%
{}
{}
{\hspace*{6pt}\titlerule\hspace*{6pt}\LARGE\bfseries \thecontentspage
	\begin{tikzpicture}[remember picture, overlay]
		\draw[fill=\mauPART, rounded corners=0pt] (2pt,0) rectangle (6,0.1pt);
	\end{tikzpicture}
}%
%~~~~~~~~~~~~~~~~~~~~~
\titlecontents{chapter}[6.5pc] % nd cách trái
{\addvspace{10pt}%
	\color{\mauCHAP}\fontsize{14pt}{1pt}\fontfamily{put}\selectfont\bfseries
}%
{\contentslabel[\chaptertitlename\,\thecontentslabel.]{6.5pc}} %nhãn
{}
{\hspace*{6pt}\titlerule\hspace*{6pt}\bfseries\thecontentspage
}%
%~~~~~~~~~~~~~~~~~~~~~
%\def\bai{\S}
\def\bai{Bài\ }
\titlecontents{section}[3.9pc]
{\addvspace{4pt}\bfseries\color{\mauSEC}}
{\fontsize{11.5pt}{13pt}\selectfont\sffamily\contentslabel[{\bai\thecontentslabel}.]{3.8pc}}
{}
{{\scriptsize \dotfill}
	\thecontentspage
}
[]
%~~~~~~~~~~~~~~~~~~~~~
\titlecontents{subsection}[6pc]
{\addvspace{1.0pt}\color{\mauSUBSEC}}
{\fontsize{13.5pt}{15pt}\selectfont\sffamily\contentslabel[\tron{\thecontentslabel}]{2.5pc}}
{}
{{\tiny\dotfill}\thecontentspage}
[]
%~~~~~~~~~~~~~~~~~~~~~
\titlecontents{subsubsection}[7.1pc]
{\addvspace{1.0pt}\color{\mauSUBSUBSEC}}
{\fontsize{12.5pt}{11pt}\selectfont\sffamily\contentslabel[\thecontentslabel.]{1pc}}
{}
{{\tiny\dotfill}\thecontentspage}
[]
%~~~~~~~~~~~~~~~~~~~~~
%--------------------------
% ĐỊNH NGHĨA CÁC MÔI TRƯỜNG 
%--------------------------
%===================================ĐNghĩa
\theoremstyle{plain} %
\theoremheaderfont{\fontfamily{put}\bfseries} %đầu
\theorembodyfont{\normalfont} % thân
\theoremseparator {.} % Ngăn cách
\newtheorem{mydn}{\color{\maudn}
	\faGg\ 
	Định nghĩa}[section]
\newtheorem{vd}{\color{\mauVD}\damEX%
	\iconVD\ Ví dụ}%[section]
\newtheorem{bt}{\color{\mauBT}\damEX\iconVD Bài}%[section]
\newtheorem{ap}{\color{\mauBT}\damEX\iconVD Áp dụng}%[section]
\newtheorem{myvidu}{\color{\mauVD}\fontfamily{put}\selectfont%
	Ví dụ}%[section]
%===================================
\theoremstyle{plain} %
\theoremheaderfont{\fontfamily{put}\bfseries} %đầu
\theorembodyfont{\slshape} % thân 
\theoremseparator {.} % Ngăn cách 
\newtheorem{mydl}{\color{\maudl}\faGg\ Định lí}%[section]
%\newtheorem{tc}{\color{\maudl}\faToggleOn\ Tính chât}%[section]
%====================================
\theoremstyle{nonumberplain} %ko đánh số, ko xuống dòng
\theoremheaderfont{\fontfamily{put}\it\bfseries} %đầu
\theorembodyfont{\normalfont} %phần thân
\theoremseparator {.} %ngăn cách
\newtheorem{nx}{\color{\maucy}\faToggleOn\,Nhận xét}
%\newtheorem{hd}{\color{\maupt}\bfseries\faCommenting\,Hướng dẫn}
\newtheorem{phantich}{\color{\maupt}\faQuestionCircle\ \sc Phân tích}
\newtheorem{tc}{\color{\maudl}\faToggleOn\ Tính chất}
%------------------------------------
\listenumerate{dn,vd,dl,tc,nx,ex,bt,vidu}%xuống dòng khi liệt kê
%\raggedbottom
%\newenvironment{tomtat}
%{\begin{tcolorbox}[
%		enhanced,breakable,
%		colback=yellow!10,
%		colframe=blue,
%		boxrule=0pt,
%		%		drop fuzzy shadow=gray,
%		left=0pt,right=0pt,top=0pt,bottom=0pt,
%		arc=0mm
%		]}
%	{\end{tcolorbox}}
\newenvironment{tomtat}{}{}
%====================================Hộp
\newtcolorbox{noidung}[1]{
	breakable, enhanced jigsaw,	frame hidden,opacityback=0,
	%	enhanced jigsaw,breakable,pad at break*=1mm,
	%	colback=pink!40,
	%	%opacityback=0, %ko nền
	%	colframe=pink!40,
	coltitle=violet,
	left=14mm,right=3mm,top=2mm,bottom=2mm,
	%--
	before skip=2mm,
	after skip=3mm,
	boxrule=0.5pt, %đường viền dày
	%toprule=1pt,
	%bottomrule=1pt,
	%leftrule=2.5pt,
	%rightrule=0pt,
	%	sharp corners,
	%	rounded corners=southwest, 
	%	rounded corners=northeast, 
	%	rounded corners=southeast,
	%	arc=7mm,
	underlay unbroken and first={
		%--khung nội dung
		\fill[fill=pink!20,drop shadow=gray] 
		([xshift=13mm]interior.north west) arc (0:-90:8mm) 
		--([xshift=5mm]interior.south west)
		--([xshift=-8mm]interior.south east) arc (-90:0:8mm) 
		--(interior.north east)
		--cycle;
		\draw[teal,line width=1pt] 
		(interior.north east) coordinate (mot)		
		--([xshift=13mm]interior.north west) coordinate (hai) arc (0:-90:8mm) coordinate (ba) 
		--([xshift=5mm]interior.south west) coordinate (bon)
		;
		\fill[teal] 
		(mot) circle (2pt)
		(hai) circle (2pt)
		(ba) circle (2pt)
		(bon) circle (2pt)
		;
	},
	underlay last={
		%--khung nội dung
		\fill[fill=pink!40,drop shadow=gray] 
		([xshift=5mm]interior.north west)
		--([xshift=5mm]interior.south west)
		--([xshift=-8mm]interior.south east) arc (-90:0:8mm) 
		--(interior.north east)
		--cycle;
		\draw[teal,line width=1pt] 
		(interior.north east) coordinate (mot)		
		--([xshift=5mm]interior.north west) coordinate (hai) 
		--([xshift=5mm]interior.south west) coordinate (bon)
		;
		\fill[teal] 
		(mot) circle (2pt)
		(hai) circle (2pt)
		(bon) circle (2pt)
		;
	},
	fonttitle=\fontfamily{qag}\bfseries\selectfont,
	title={#1\phantom{T}},
	attach boxed title to top left={xshift=1cm,yshift=-1mm},
	boxed title style={empty},
	%	underlay boxed title={
	underlay unbroken and first={
		\fill[violet!70] ([xshift=5mm]interior.north west) node[white,scale=1.5]{
			%\faBell
			\faEdit
		} circle (6.5mm) ;
		\draw[white] ([xshift=5mm]interior.north west) circle (6mm) ;
	},
	%----
	%fontupper=\sffamily % chữ nghiêng phần nội dung
}
%--------------------------
%\def\maudl{orange}\def\nendl{yellow!5}
\newenvironment{boxdl}
{\begin{tcolorbox}
		[enhanced jigsaw,breakable,pad at break*=1mm,
		colback=\nendl,boxrule=0pt,frame hidden, 
		opacityback=0.5,
		left=2mm, right=0pt, bottom=1.5pt, top=1.5pt,
		before skip=2mm,
		after skip=2mm,
		%		borderline west={1mm}{0cm}{green!70!black}
		overlay={\draw[double,line width=1.5pt,\maudl] ([xshift=3pt]interior.north west)--([xshift=3pt]interior.south west);}
		]}%0mm lề trái
	{\end{tcolorbox}}
%\def\maudn{green!70!black}\def\nendn{green!10}
\newenvironment{boxdn}
{\begin{tcolorbox}
		[enhanced jigsaw,breakable,pad at break*=1mm,
		%		drop fuzzy shadow=orange,
		colback=\nendn,boxrule=0pt,frame hidden, 
		opacityback=0.5,
		left=2mm, right=0pt, bottom=1.5pt, top=1.5pt,
		before skip=2mm,
		after skip=2mm,
		%		borderline west={1mm}{0cm}{orange}
		overlay={\draw[double,line width=1.5pt,\maudn] ([xshift=3pt]interior.north west)--([xshift=3pt]interior.south west);}
		]}%0mm lề trái
	{\end{tcolorbox}}

\newenvironment{boxkn}
{\begin{tcolorbox}
		[enhanced jigsaw,breakable,pad at break*=1mm,
		%opacityback=0,
		boxrule=0pt,frame hidden, 
		colback=white,
		opacityback=0,
		left=9mm, right=0pt, bottom=1pt, top=1pt,
		before skip=2mm,
		after skip=2mm,
		toprule=1pt,
		borderline west={1mm}{0.39cm}{\maucy},
		underlay unbroken and first={
			\fill[white,draw=\maucy,line width=1pt] ([xshift=0.4cm,yshift=0cm]interior.west) circle (0.4cm);
			\draw ([xshift=0.4cm,yshift=0cm]interior.west) node[\maucy]{\large\bfseries \faKey};
		},
		]}
	{\end{tcolorbox}}



%-- Khung vidu
\newenvironment{boxvidu}
{\begin{tcolorbox}
		[enhanced jigsaw,breakable,pad at break*=1mm,
		%		drop fuzzy shadow=orange,
		colback=white,boxrule=0pt,frame hidden, 
		opacityback=0,
		left=2mm, right=0pt, bottom=1.5pt, top=1.5pt,
		before skip=2mm,
		after skip=2mm,
		%		borderline west={1mm}{0cm}{orange}
		overlay={\draw[line width=1pt,\mauVD!70] ([xshift=2pt]interior.north west)--([xshift=2pt]interior.south west);}
		]}%0mm lề trái
	{\end{tcolorbox}}
\newenvironment{dn}{\begin{boxdn}\begin{mydn}}{\end{mydn}\end{boxdn}}
\newenvironment{dl}{\begin{boxdl}\begin{mydl}}{\end{mydl}\end{boxdl}}
%--------------------Chú ý
\newenvironment{note}
{\begin{tcolorbox}
		[enhanced jigsaw,breakable,pad at break*=1mm,
		opacityback=0,boxrule=0pt,frame hidden,
		left=8mm, right=0pt, bottom=0pt, top=0pt,
		before skip=1mm,
		after skip=1mm,
		underlay unbroken and first={
			\draw ([xshift=0.3cm,yshift=-0.32cm]interior.north west) node[\maucy]{\large\bfseries \faExclamationTriangle};
		},
		fontupper=\it,
		]}
	{\end{tcolorbox}}
\setlength{\parindent}{0pt} %không thụt đầu dòng
%--------------------Chú ý
\newenvironment{luuy}
{\begin{tcolorbox}
		[enhanced jigsaw,breakable,pad at break*=1mm,
		opacityback=0,boxrule=0pt,frame hidden,
		left=8mm, right=0pt, bottom=0pt, top=0pt,
		before skip=1mm,
		after skip=1mm,
		underlay unbroken and first={
			\draw ([xshift=0.3cm,yshift=-0.32cm]interior.north west) node[\maucy]{\large\bfseries \faExclamationTriangle};
		},
		fontupper=\it,
		]}
	{\end{tcolorbox}}
\setlength{\parindent}{0pt} %không thụt đầu dòng
%==================
%--------------------đáp số
\renewcommand{\dapso}[1]{\hspace{\fill}\mbox{}\linebreak[0]\hspace*{\fill}\mbox{\fontsize{8pt}{1pt}\selectfont\color{\maupt}{\color{\maucy}{\faKey}}~#1}
}
%====================
\usepackage[final]{pdfpages} % input file bìa pdf
%----Cài đặt \boxEX -ô vuông sd tài liệu
\def\hideboxEX{}
\renewcommand{\boxEX}[2][15pt]{%
	\setbox1=\hbox{#2\strut}
	\def\sizeboxEX{#1}%Thay #1 bằng 15pt để cố định tất cả boxEX
	\ \EXbox{\hideboxEX#2\strut}\ %
}
%\renewcommand{\boxEX}[2][15pt]{%
%	\setbox1=\hbox{#2}%
%	\def\sizeboxEX{#1}%
%	\ \EXbox{#2}\ %
%}
\def\AnboxEX{
	\renewcommand{\boxEX}[2][15pt]{%
		\def\sizeboxEX{##1}%
		\setbox1=\hbox{##2}%
		\ \EXbox{}\ %
}}
%===Định nghĩa in Mục lục chính
\def\muclucchinh{
	\thispagestyle{empty}
	\pagenumbering{roman}%đánh số trang dạng i,ii,...
	\tableofcontents %lệnh in mục lục chính
	%	\cleardoublepage %tạo trang trống
}
%-------------Môi trường thủ công 2 cột Dotted
\newcommand{\dongcham}[1]{
	\def\sod{#1}
	\pgfmathsetmacro{\sodong}{2*\sod -1} 
	\columnsep=10pt
	\vspace*{-3.5mm}
	\begin{multicols}{2}
		\foreach \dotline in{1,...,\sodong}
		{\noindent\color{\mauDC}{\dotfill}\\[1mm]
		}\noindent\color{\mauDC}{\dotfill}\\[-4mm]
	\end{multicols}
}
%--------------Cài đặt lại dòng kẻ \dotline
%---ĐN dòng chấm kiểu 1 cột
\def\mauDC{\mauLG!50}
\def\dotlineOne{
	\pgfmathsetmacro{\sodong}{\numlinedot-1} 
	\par
	\foreach \dotline in{1,...,\sodong}
	{\noindent{\color{\mauDC}\dotfill}\\[1mm]
	}{\noindent{\color{\mauDC}\dotfill}}\\[-2mm]
}
%--SD các kiểu đn trên vào lệnh dotlineEX
\renewcommand{\dotlineEX}[1]{
	\setlength{\parindent}{0pt}%
	\def\numlinedot{#1}
	\dotlineOne % dòng chấm bình thường (1 cột)
}
%--------------- Đặt lại ĐN khung nền sau khi sử dụng Lệnh \tieude{}
%\AfterEndEnvironment{ex}{
%	\renewcommand{\dnsub}[1]{
%		\begin{tikzpicture}[line join=round,line cap=round]%
%			\node[inner sep=3pt,text=white,rounded corners=0pt,fill=\mausubsec] (muccon) at (0,0){\hspace*{2mm}#1\hspace*{0.75cm}}; %trái
%			\fill[fill=orange] ([xshift=-3pt]muccon.north east) --([xshift=-3pt]muccon.east) 
%			.. controls +(-90:0.2) and +(10:0.2) ..  ([xshift=-0.35cm]muccon.south east)
%			--(muccon.south east) .. controls +(10:0.2) and +(-90:0.2) .. ([xshift=0.35cm]muccon.east)--([xshift=0.35cm]muccon.north east)--cycle ; %phải
%		\end{tikzpicture}%
%	}
%}
%=================ĐN KHUNG
\def\endbox{\end{tcolorbox}}
%--khung ex
\def\beginboxEX{\begin{tcolorbox}[%
	enhanced jigsaw,breakable,colback=\nenEX,
	boxrule=0.6pt,arc=0pt,
	opacityback=0.5,
	left=2mm,right=2mm,top=2mm,bottom=2mm,
	before skip=2mm,after skip=3mm,
	colframe=\mauEX!80!black,drop lifted shadow=\mauEX]}
%--Khung bt
\def\beginboxBT{\begin{tcolorbox}[%
		enhanced jigsaw,breakable,colback=\nenBT,
		boxrule=0.6pt,arc=0pt,
		opacityback=0.5,
		left=2mm,right=2mm,top=2mm,bottom=2mm,
		before skip=2mm,after skip=3mm,
		colframe=\mauBT!80!black,drop lifted shadow=\mauBT]}
	%-- Khung vd
	\def\beginboxVD{
		\begin{tcolorbox}[
			enhanced jigsaw,breakable,
			arc=0mm,colback=\nenVD,
			boxrule=0.6pt,arc=0pt,
			opacityback=0.5,
			colframe=\mauVD!80!black,drop lifted shadow=\mauVD,
			left=2mm,right=2mm,top=2mm,bottom=2mm,
			before skip=2mm,after skip=3mm,
			]
		}
		%---
		%=======Đn các phương án
		%		\def\mauDA{\mauEX}
		%		\def\mauTrue{red}
		\def\khoanhtrondapan{
			\renewcommand*\circled[1]{\tikz[baseline=(char.base)]{
					\node[shape=circle,draw=\mauDA,inner sep=1pt] (char) {##1};}}
			\renewcommand{\TrueEX}{\stepcounter{dapan}
				{\squareEX{\textbf{\fontfamily{qag}\selectfont\color{\mauDA}\Alph{dapan}}}} \ignorespaces}
			\renewcommand{\FalseEX}{\stepcounter{dapan}
				{\circled{\textbf{\fontfamily{qag}\selectfont\color{\mauDA}\Alph{dapan}}}} \ignorespaces}
			%---Chọn đáp án
			\renewcommand*\circEX[1]{\tikz[baseline=(char.base)]{\node[shape=circle,draw=\mauTrue,inner sep=1pt,fill=\mauTrue!3] (char) {\color{\mauTrue} ##1};}}
		}
		%----
		\def\khongkhoanhtrondapan{

			\renewcommand{\TrueEX}{\stepcounter{dapan}
				{\squareEX{\textbf{\fontfamily{qag}\selectfont\color{\mauDA}\Alph{dapan}}}} \ignorespaces}
			\renewcommand{\FalseEX}{\stepcounter{dapan}
				{\textbf{\fontfamily{qag}\selectfont\color{\mauDA}\Alph{dapan}.}} \ignorespaces}
			%---Chọn đáp án

			\renewcommand*\circEX[1]{\tikz[baseline=(char.base)]{\node[shape=circle,draw=\mauTrue,inner sep=1pt,fill=\mauTrue!3] (char) {\color{\mauTrue} ##1};}}
		}
		%%--------------------------Đn lại màu sắc theo từng môi trường
		%\AtBeginEnvironment{bt}{\def\mauDA{\mauBT}}
		%\AtBeginEnvironment{ex}{\def\mauDA{\mauEX}}
		%\AtBeginEnvironment{vd}{\def\mauDA{\mauVD}}
		%===========================MainDeThi
		%FULL WIDTH
		\def\FULLWIDTH{
			\newpage
			\fancyhead[LO,RE]{\headmucluc}
			\def\headchuong{\headmucluc}
			\newgeometry{top=\tren cm, bottom=\duoi cm, left=\trai cm, right=\phai cm}
			\def\mepphai{\phai}
		}
		\def\NOTE{
			\newpage
			\fancyhead[LO,RE]{\myfancyhead}
			\def\headchuong{\myfancyhead}
			\pgfmathsetmacro{\mepphai}{\phai+\leftnote}
			\newgeometry{top=\tren cm, bottom=\duoi cm, left=\trai cm, right=\mepphai cm}
		}
		%==================
		\def\BTRL{}
	%======================================================== MÔI TRƯỜNG VIDUII
	\colorlet{stardrawcol}{\mausao}
	\colorlet{starmkdcol}{\mausao}
	\colorlet{starempcol}{white}

	\newcommand\score[2]{%
		\pgfmathsetmacro\pgfxa{#1 + 1}%
		\tikzstyle{scorestars}=[star, star points=5, star point ratio=2, draw=stardrawcol, inner sep=2pt, anchor=outer point 3]%

		\begin{tikzpicture}[baseline]
			\foreach \i in {1, ..., #2} {
				\pgfmathparse{\i<=#1 ? "starmkdcol" : "starempcol"}
				\edef\starcolor{\pgfmathresult}
				\draw (\i*3.5ex, 0ex) node[name=star\i, scorestars, fill=\starcolor]  {};
			}
		\end{tikzpicture}%
	}
	\newcommand{\mkstar}[1]{\protect\score{#1}{4}}
% ====================================================
% lệnh siunit 
\newcommand{\xsi}[2]{\SI[parse-numbers=false]{#1}{#2}}
%====================================================
% --- ký hiệu suất điện động và công suất như sgk
% ký hiệu từ font Boondox 
\DeclareFontFamily{U}{BOONDOX-cal}{\skewchar\font=45 }
\DeclareFontShape{U}{BOONDOX-cal}{m}{n}{
	<-> s*[1.05] BOONDOX-r-cal}{}
\DeclareFontShape{U}{BOONDOX-cal}{b}{n}{
	<-> s*[1.05] BOONDOX-b-cal}{}
\DeclareMathAlphabet{\bdx}{U}{BOONDOX-cal}{m}{n}
\SetMathAlphabet{\bdx}{bold}{U}{BOONDOX-cal}{b}{n}
\DeclareMathAlphabet{\bbdx}{U}{BOONDOX-cal}{b}{n}
\newcommand{\calE}{\bdx{E}}
\newcommand{\calP}{\bdx{P}}